%%
%% The abstract is a short summary of the work to be presented in the
%% article.
\begin{abstract}
Scanned document images frequently suffer from noise, uneven illumination, and low contrast, all of which degrade the accuracy of downstream Optical Character Recognition (OCR) systems. This paper presents a high-performance, modular C++ pipeline for document image binarization that addresses these challenges through a configurable sequence of preprocessing, binarization, and postprocessing stages. The preprocessing chain combines median denoising, separable box-blur background estimation and removal, and percentile-based contrast stretching. Five binarization algorithms are provided: the global Otsu method, and the local adaptive methods of Sauvola, NICK, Su-Lu-Tan, and a proposed adaptive technique that jointly exploits global image statistics and normalized local standard deviation. Postprocessing includes morphological opening and closing operations as well as border cleanup. Computational efficiency is achieved through integral images, which reduce the complexity of local-statistics computation to constant time per pixel, and through OpenMP parallelization applied across all compute-intensive stages. The pipeline is validated with a Catch2-based test suite comprising 56 test cases and 1038 assertions. A qualitative evaluation on a degraded Arabic calligraphy document demonstrates that preprocessing is essential for local adaptive methods and that the choice of binarization algorithm significantly affects output quality.
\end{abstract}
